\documentclass[12pt,a4paper]{article}
\usepackage[margin=1in]{geometry}
\usepackage{amsmath,amssymb,amsthm}
\usepackage{graphicx}
\usepackage{url}
\usepackage{cite}
\usepackage{xcolor}
\definecolor{navyDark}{HTML}{003f88}
\definecolor{navy}{HTML}{0066cc}
\usepackage[colorlinks=true,allcolors=navyDark,citecolor=navy]{hyperref}

\usepackage{algorithm}
\usepackage{algorithmicx, algpseudocode}
\usepackage{tikz}
\usepackage{pgfplots}
\pgfplotsset{compat=1.18}

% Custom theorem environments
\newtheorem{theorem}{Theorem}
\newtheorem{lemma}{Lemma}
\newtheorem{definition}{Definition}
\newtheorem{proposition}{Proposition}
\newtheorem{corollary}{Corollary}

\renewcommand{\algorithmicrequire}{\textbf{Input:}}
\renewcommand{\algorithmicensure}{\textbf{Output:}}

% Custom commands for common notation
\newcommand{\F}{\ensuremath{\mathbb{F}} }
\newcommand{\Fq}{\ensuremath{\F_q} }
\newcommand{\Fe}[1]{\ensuremath{\F_{q^{#1}}} } % Field extension
\newcommand{\G}{\ensuremath{\mathbb{G}} }
\newcommand{\gComm}[1]{\color{gray}\Comment{#1}\color{black}}
\newcommand{\gText}[1]{\color{gray}#1\color{black}}

\newcommand{\IOP}[4]{\fbox{\parbox{\linewidth}{
\begin{center}
    \begin{tikzpicture}[node distance=13cm, every node/.style={font=\footnotesize}]
    \node (L)[draw,font=\large,rectangle] {#2};\node (R)[draw,font=\large,rectangle, right of=L] {#3};
    #4
    \end{tikzpicture}\linebreak[1]

    #1
\end{center}
}}
}

\newcommand{\TikzArrow}[4]{
    \draw[->] ([yshift=-{#4}cm]#1.center) -- ([yshift=-{#4}cm]#2.center) node[midway, above] {#3};
}

\newcommand{\LDoes}[2]{\node[above=-2pt, right] at ([yshift=-{#2}cm,xshift=-.7cm]L.center) {#1};}
\newcommand{\RDoes}[2]{\node[above=-2pt, left] at ([yshift=-{#2}cm,xshift=.7cm]R.center) {#1};}
\newcommand{\RtoL}[2]{\TikzArrow{R}{L}{#1}{#2}}
\newcommand{\LtoR}[2]{\TikzArrow{L}{R}{#1}{#2}}

\title{Probabilistic Proof for Pairings}

\author{Shramee Srivastav \\ 
MIST.cash|FOCBB, Shhtarknet \\ 
Email: shramee@proton.me}

\date{\today}

\emergencystretch=1em
\begin{document}

\maketitle

\begin{abstract}
Elliptic curve pairings are fundamental to modern zero-knowledge proof systems, including Groth16, PLONK, and KZG polynomial commitments. In resource-constrained execution environments (Ethereum VM, BitVM, blockchain smart contracts) where computation is metered by gas fees, and in arithmetized computation systems (R1CS, Plonkish, Cairo AIR) where each computational step imposes significant prover overhead, the cost of pairing computation remains a practical bottleneck.

We present a public-coin interactive proof system that consolidates the Miller loop iterations and final exponentiation into a single polynomial relation. The relation can then be verified as polynomial identities greatly reducing the pairing costs.

We present a full implementation in gnark~\cite{gnark}---the industry-standard framework for arithmetized circuit development---as a reusable polynomial ring toolkit in the \texttt{emulated} package. Benchmarked on a Groth16 verification simulation (2 fixed-point pairings + 1 full pairing + 1 $\mathbb{F}_{p^{12}}$ multiplication) on BN254, our gnark implementation achieves a \textbf{39\% reduction in SCS constraints} (1,890,771~$\to$~1,158,194), a \textbf{60\% reduction in commitments} (739,198~$\to$~295,946), and a $\sim$\textbf{42\%} reduction in compile-time memory and up to \textbf{52\%} in solver memory. These improvements hold across both SCS and R1CS constraint systems and enable practical pairing verification in recursive proof systems and resource-constrained environments.
\end{abstract}

\tableofcontents
\listofalgorithms
\listoftables

\section{Introduction}

Blockchain technology continues to mature with improving scalability and infrastructure. However, privacy remains a significant barrier to entry for many users and applications. Zero-knowledge proofs (ZKPs) have emerged as a powerful solution to this challenge.

\subsection{Motivation}


Elliptic curve pairings represent fundamental building blocks for numerous cryptographic systems including Groth16\cite{Gro16}, PLONK\cite{plnk19}, BLS signature schemes\cite{bls05}, and KZG polynomial commitment schemes\cite{kzg10}.

However, \textbf{pairings are computationally intensive}. This creates critical bottlenecks in two important deployment contexts:

\begin{enumerate}
\item \textbf{Resource-constrained virtual machines} (Ethereum VM, BitVM, ZKVMs): Computation is metered — each operation incurs gas costs. A single pairing verification requiring thousands of field operations becomes economically prohibitive especially during peak network congestion.

\item \textbf{Arithmatized computation systems} (R1CS~\cite{r1cs13}, Plonkish~\cite{plnk19}, Cairo AIR~\cite{cairo21}): Each operation must be represented algebraically, adding significant prover overhead beyond native execution costs.
\end{enumerate}

In such contexts, we can make pairing verification practical by using an efficient proof system to verify the correctness of the pairing computation, rather than executing all operations natively within the expensive environment.

\subsection{Our Contribution}

We present a public-coin interactive proof system for verifying pairing computations that consolidates verification into unified polynomial relationships under the Schwartz-Zippel lemma. By treating complete Miller loop iterations as single polynomial relationships rather than sequences of individual field operations\cite{FXFM23}, our approach achieves:

\begin{itemize}
\item \textbf{45\% reduction} in Miller loop field operations
\item \textbf{77\% reduction} in commitment overhead (Fiat-Shamir hashing)
\item \textbf{37\% overall speedup} for complete pairing verification
\end{itemize}

Our key insight is that instead of verifying individual field operations separately, we can verify entire Miller loop iterations as single polynomial identities. This consolidation dramatically reduces both the number of modular operations required and the communication overhead from intermediate commitments.

\section{Background}

\begin{algorithm}
\caption{Optimal Ate pairing over Barreto-Naehrig curves~\cite{hsp10}}
\begin{algorithmic}[1]
\Require{$P \in \G_1, Q \in \G_2$}
\Ensure{$a_{opt}(Q, P)$}
\State{Write $s = 6t + 2 = \sum_{i=0}^{L-1} s_i 2^i$ where $s_i \in \{-1, 0, 1\}$ and $s_L = 1$}
\State{$T \leftarrow Q, f \leftarrow 1$}
\For{$i = L - 2$ to $0$}
    \State{$f \leftarrow f^2 \cdot l_{T,T}(P)$}
    \State{$T \leftarrow [2]T$}
    \If{$s_i = 1$}
        \State{$f \leftarrow f \cdot l_{T,Q}(P)$}
        \State{$T \leftarrow T + Q$}
    \EndIf
    \If{$s_i = -1$}
        \State{$f \leftarrow f \cdot l_{T,-Q}(P)$}
        \State{$T \leftarrow T - Q$}
    \EndIf
\EndFor
\State{$Q_1 \leftarrow \pi_p(Q), Q_2 \leftarrow \pi_p^2(Q)$}
\State{$f \leftarrow f \cdot l_{T,Q_1}(P), T \leftarrow T + Q_1$}
\State{$f \leftarrow f \cdot l_{T,-Q_2}(P), T \leftarrow T - Q_2$}
\State{$f \leftarrow f^{p^{12} - 1 / r}$}
\State{\Return $f$}
\end{algorithmic}
\end{algorithm}

\subsection{Pairing Computation}

An elliptic curve pairing is a non-degenerate bilinear map $e: \G_1 \times \G_2 \rightarrow \G_T$, where $\G_1$ and $\G_2$ are groups of points on elliptic curves over finite fields, and $\G_T$ is a multiplicative group in an extension field. For practical cryptographic applications, we require that this map be efficiently computable. We will focus on the optimal ate pairing~\cite{hsp10} over Barreto-Naehrig (BN) curves, which is widely used due to its efficiency. We will specifically use the Ethereum's BN254 curve, defined over a 254-bit prime field $\Fq$, with groups $\G_1, \G_2$ of order $r$ and target group $\G_T$ contained in the 12th extension field $\Fe{12}$.

The computation of an optimal ate pairing---the focus of our implementation---consists of two primary phases:

\begin{enumerate}
    \item \textbf{Miller loop}: The core iterative computation that evaluates line functions at pairs of points from $\G_1$ and $\G_2$, accumulating intermediate results to produce a value in $\Fe{12}$.
    
    \item \textbf{Final Exponentiation}: This step raises the Miller loop output to the power $(q^{12}-1)/r$, where $r$ is the order of the groups, ensuring the result lies in the correct subgroup $\G_T$.
\end{enumerate}

While both phases present optimization opportunities, the constraint-based verification of these computations introduces unique challenges and possibilities that differ significantly from native implementations.

\subsection{Schwartz-Zippel lemma} \label{sec:bg:sz}

The Schwartz-Zippel lemma (or more precisely, the Demillo-Lipton-Schwartz-Zippel lemma~\cite{DemLip78,Sch79,Zip79,Sch80}) provides a probabilistic method for testing polynomial identities. The lemma states that for a non-zero polynomial $P(x_1, x_2, \ldots, x_n)$ of total degree $d$ over a finite field $\F$, when evaluated at uniformly random field elements, the probability that $P$ evaluates to zero is at most $d/|\F|$.

This probabilistic bound becomes negligible for cryptographically sized fields where $d \ll |\F|$, making it practical to verify polynomial identities by evaluation instead of coefficient-wise operations.

\subsection{Constraint-Based Computation Models}

\subsubsection{Computational Model}

Our optimization targets constraint-based computation models where programs are represented as systems of algebraic constraints over finite fields. We formalize this model as follows:

\textbf{Definition (Constraint System):} A constraint system over a finite field $\Fq$ consists of:
\begin{itemize}
\item A set of variables $\mathbf{v} = (v_1, \ldots, v_m) \in \Fq^m$
\item A partition of variables into:
\begin{itemize}
\item \textit{Public inputs} $\mathbf{x} \in \Fq^n$ (known to verifier)
\item \textit{Private witnesses} $\mathbf{w} \in \Fq^{m-n}$ (known only to prover)
\end{itemize}
\item A set of polynomial constraints $\{C_i(\mathbf{v}) = 0\}_{i=1}^{\ell}$ where each $C_i$ is a polynomial over $\Fq$
\end{itemize}

A valid computation satisfies all constraints when variables are assigned appropriate values.

\textbf{Witnesses and Hints:} In constraint systems, complex operations (e.g., field division, polynomial quotients) can be computed efficiently using \textit{witnesses} or \textit{hints}---auxiliary values provided by the prover but verified by constraints. For instance, to verify $a/b = c$:
\begin{enumerate}
\item Prover provides $c$ as a witness
\item Verifier checks the constraint $b \cdot c = a$
\end{enumerate}

This approach transforms expensive operations into cheap verifications, fundamentally altering the cost model.

\textbf{Examples of Constraint Systems:}
\begin{itemize}
\item \textbf{R1CS} (Rank-1 Constraint Systems)\cite{r1cs13}: Constraints of the form $(A\mathbf{v}) \circ (B\mathbf{v}) = (C\mathbf{v})$ where $\circ$ is element-wise multiplication
\item \textbf{Plonkish}\cite{plnk19}: Gate-based constraints with custom gates and copy constraints
\item \textbf{Cairo/AIR}\cite{cairo}: Algebraic Intermediate Representation with polynomial constraints over execution traces
\end{itemize}

\subsubsection{Cost Model}

We measure computational cost in terms of \textit{constraints} (denoted by \textbf{C}), representing the fundamental unit of work in constraint systems. One constraint typically corresponds to one multiplication gate or one algebraic relation to verify.

\textbf{Basic Field Operations in $\Fq$:}
\begin{itemize}
\item \textbf{Addition/Subtraction:} 0C (constraints are linear, free to compose)
\item \textbf{Multiplication:} 1C (one constraint)
\item \textbf{Constant multiplication:} 0C (absorbed into constraint coefficients)
\item \textbf{Division/Inversion:} 2C (one multiplication constraint + one equality check)
\end{itemize}

\textbf{Extension Field Operations:} For $\Fe{k}$ elements represented as $k$ coefficients in $\Fq$:
\begin{itemize}
\item \textbf{Addition:} 0C (component-wise, linear)
\item \textbf{Multiplication:} Traditional approach requires $O(k^2)$ base field multiplications, but our polynomial verification reduces this significantly
\item \textbf{Evaluation at point $x$:} $k \cdot 1C = kC$ (Horner's method: $k-1$ multiplications + $k-1$ additions)
\end{itemize}

\textbf{Commitment Overhead:} Beyond constraint costs, we account for \textit{calldata} or \textit{commitment overhead}---the number of field elements the prover must commit to (via Fiat-Shamir hashing in non-interactive proofs). For blockchain deployments, this overhead directly impacts transaction costs and is often the dominant factor.

\subsubsection{Cost Model Comparison}

Table \ref{tab:cost_comparison} contrasts costs between native implementations and constraint systems:

\begin{table}[h]
\centering
\caption{Cost comparison: Native vs Constraint-based computation}
\label{tab:cost_comparison}
\begin{tabular}{|l|c|c|c|}
\hline
\textbf{Operation} & \textbf{Native Cost} & \textbf{Constraint Cost} & \textbf{Ratio} \\
\hline
$\Fq$ Multiplication & 1$\times$ & 1C & 1$\times$ \\
$\Fq$ Inversion & 100$\times$ & 2C & $\approx$50$\times$ cheaper \\
$\Fe{12}$ Multiplication & 50$\times$ & 54C (traditional) & Similar \\
$\Fe{12}$ Multiplication & 50$\times$ & 39C (our method) & Better \\
\hline
\end{tabular}
\end{table}

\textbf{Key Insight:} Operations that are traditionally expensive (inversions, divisions) become cheap in constraint systems when computed via witnesses. This inverts traditional optimization priorities: we can afford inversions to save multiplications, enabling affine coordinate systems and other techniques impractical in native implementations.

Recognizing and exploiting these unique computational properties is essential for unlocking hidden performance potential. Youssef El Housni's pioneering work \textbf{Pairings in Rank-1 Constraint Systems}~\cite{PR1CS22} represents, to our knowledge, the first systematic exploration of these optimization opportunities, providing a breakthrough in the practicality of recursive proving by thoughtfully leveraging the distinctive cost model inherent to constraint-based systems.
\subsection{Related Work}

The landscape of efficient pairing verification in constraint systems has evolved through several key contributions. El Housni's \textbf{Pairings in Rank-1 Constraint Systems}~\cite{PR1CS22} established the theoretical framework for optimizing pairings in Rank-1 Constraint Systems, introducing techniques for efficient Miller loop implementation and final exponentiation that exploit the unique properties of constraint-based computation models. This work demonstrated that operations traditionally considered expensive in native implementations could be restructured to achieve remarkable efficiency gains in arithmetized environments.

Feltroidprime's subsequent research on \textbf{enhanced extension field multiplication techniques}~\cite{FXFM23} introduced polynomial verification methods using the Schwartz-Zippel lemma for $\Fe{12}$ arithmetic operations. By representing extension field elements as polynomials and verifying their operations through polynomial arithmetic, this approach achieved significant performance improvements in the pairing verification.

Novakovic and Eagen's work \textbf{On Proving Pairings}~\cite{OPP24} presented breakthrough optimization strategies that enable the elimination of final exponentiation overhead through residue witness techniques. Their approach demonstrates that two elements in $\Fe{12}$ can be proven equivalent without expensive final exponentiation by introducing witness elements that satisfy specific algebraic relationships, and effectively embedding the final exponentiation verification into the Miller loop computation itself.

\newcommand{\InlineFor}[2]{\State \textbf{for} #1 \textbf{do} #2}

\section{Pairing Operations as Polynomials}

Algorithm~\ref{alg:pairing} gives the pairing computation combining El Housni's R1CS optimizations~\cite{PR1CS22} with Novakovic-Eagen's residue witness technique~\cite{OPP24}. The 27th root of unity $W$ scales the Miller loop output $f$ to be a cubic residue (see~\cite{OPP24} §4.3); if $f$ is not already a cubic residue, multiplying by $W^w$, $w\in\{1,2\}$ corrects this.

Our contribution is to express each step of this algorithm as a single multi-operand polynomial product verified via the Schwartz-Zippel lemma, rather than as a sequence of individual binary multiplications as in~\cite{FXFM23}.

\begin{algorithm}
    \caption{Optimal Ate Multi-Pairing with Residue witness check}
    \label{alg:multipairing}
    \begin{algorithmic}[1]
    \Require Vectors $\vec{P} = (P_1, \dots, P_N) \in \G_1^N$, $\vec{Q} = (Q_1, \dots, Q_N) \in \G_2^N$, $c \in \Fe{12}, w \in \{0,1,2\}$
    \Statex $\text{Internal } W, W^2 \in \Fe{12}$ \Comment{$W$ is 27th root of unity}
    \Ensure $\prod_{k=1}^N a_{opt}(Q_k, P_k)$
    \State Write $s = 6t + 2 = \sum_{i=0}^{L-1} s_i 2^i$ where $s_i \in \{-1, 0, 1\}$ and $s_L = 1$
    \State $c^{-1} \leftarrow 1 / c$ \Comment{Inverse Residue witness}
    \State $f \leftarrow c^{-1}$ \Comment{Initialize with inverse residue witness}

    \InlineFor {$k = 1$ to $N$} {$T_k \leftarrow Q_k$}

    \Statex

    \For{$i = L - 2$ to $0$}
        \State $f \leftarrow f^2$ \label{pr:ml:sq}
        \State $f \leftarrow f \cdot \prod_{k=1}^N l_{T_k,T_k}(P_k)$ \label{pr:ml:ldbl}

        \InlineFor {$k = 1$ to $N$} {$T_k \leftarrow [2]T_k$}

        \If{$s_i = 1$}
            \State $f \leftarrow f \cdot c^{-1}$ \Comment{Residue witness}\label{pr:ml:rw1bit}
            \State $f \leftarrow f \cdot \prod_{k=1}^N l_{T_k,Q_k}(P_k)$ \label{pr:ml:ladd1bit}
            \InlineFor {$k = 1$ to $N$} {$T_k \leftarrow T_k + Q_k$}

        \ElsIf{$s_i = -1$}
            \State $f \leftarrow f \cdot c$ \Comment{Residue witness}\label{pr:ml:rw-1bit}
            \State $f \leftarrow f \cdot \prod_{k=1}^N l_{T_k,-Q_k}(P_k)$ \label{pr:ml:ladd-1bit}
            \InlineFor {$k = 1$ to $N$} {$T_k \leftarrow T_k - Q_k$}

        \EndIf
    \EndFor
    \Statex

    \InlineFor {$k = 1$ to $N$} {$Q_{1,k} \leftarrow \pi_p(Q_k)$, $Q_{2,k} \leftarrow \pi_p^2(Q_k)$}

    \State $f \leftarrow f \cdot \prod_{k=1}^N l_{T_k,Q_{1,k}}(P_k)$ \label{pr:ml:pi1}

    \InlineFor {$k = 1$ to $N$} {$T_k \leftarrow T_k + Q_{1,k}$}

    \State $f \leftarrow f \cdot \prod_{k=1}^N l_{T_k,-Q_{2,k}}(P_k)$ \label{pr:ml:pi2}

    \Statex
    \If{$w = 1$}
        \State $f \leftarrow f \cdot W$ \label{pr:ml:27r}
    \ElsIf{$w = 2$}
        \State $f \leftarrow f \cdot W^2$ \label{pr:ml:27r2}
    \EndIf

    \State $f \leftarrow f \cdot c^{-q} \cdot c^{q^2} \cdot c^{-q^3}$ \Comment{Uses Frobenius maps, Negative exponents use $c^{-1}$}\label{pr:ml:fe}

    \Statex

    \State \Return $f$
    \end{algorithmic}
\end{algorithm}

\subsection{Comparative Analysis~\label{sec:ops:comp-analysis}}

We compare the two schemes for a 3-pair Miller loop zero-bit step (the most common case in BN254's NAF\cite{naf} representation of $s = 6t+2$). Here $F \in \Fe{12}$ is the accumulator, $L_i \in \text{Sparse}_{01379}$ are the three doubling line functions, $R \in \Fe{12}$ is the updated accumulator, and $Q$ is the quotient polynomial. $\text{Sparse}_{01379}$ denotes the subspace of $\Fe{12}$ elements whose only nonzero coefficients are at degrees $0,1,3,7,9$---the form that BN254 line functions naturally take~\cite{PR1CS22}, reducing their evaluation cost from 12C to 4C.

Adapting~\cite{FXFM23} to lines \ref{pr:ml:sq}--\ref{pr:ml:ldbl} of Algorithm~\ref{alg:pairing} chains four binary multiplications, each producing an intermediate remainder that must be committed:

\begin{align*}
    F(x) \cdot F(x) &\stackrel{?}{=} Q_1(x) \cdot P_{12}(x) + T_1(x) \\
    T_1(x) \cdot L_1(x) &\stackrel{?}{=} Q_2(x) \cdot P_{12}(x) + T_2(x) \\
    T_2(x) \cdot L_2(x) &\stackrel{?}{=} Q_3(x) \cdot P_{12}(x) + T_3(x) \\
    T_3(x) \cdot L_3(x) &\stackrel{?}{=} Q(x) \cdot P_{12}(x) + R(x)
\end{align*}

Each of the five polynomial evaluations ($F$, $T_1$--$T_3$, $R$) costs 12C and each line costs 4C, with all four remainders requiring commitment: \textbf{76C and 48 $\Fq$ commitments}.

Our scheme collapses this into a single identity:\label{sec:op:comp:our}

$$F(x)^2 \cdot L_1(x) \cdot L_2(x) \cdot L_3(x) \stackrel{?}{=} Q(x) \cdot P_{12}(x) + R(x)$$

Only $R$ requires commitment, and only $F$ and $R$ need full 12C evaluation; the three sparse lines cost 4C each: \textbf{41C and 12 $\Fq$ commitments} --- a \textbf{46\% constraint reduction} and \textbf{75\% fewer commitments}. The comparison across all three schemes is summarised in Table~\ref{tab:miller_loop_comparison}.

\begin{table}[h]
\centering
\caption{3-pair Miller loop zero-bit step comparison}
\label{tab:miller_loop_comparison}
\begin{tabular}{|l|c|c|c|}
\hline
\textbf{Scheme} & \textbf{Constraints} & \textbf{$\Fq$ Commitments} & \textbf{vs baseline} \\
\hline
\cite{PR1CS22} & 132C & --- & --- \\
\hline
\cite{FXFM23} & 76C & 48 & $-$42\% \\
\hline
\textbf{Ours} & \textbf{41C} & \textbf{12} & $-$\textbf{69\%} \\
\hline
\multicolumn{4}{|l|}{\footnotesize \cite{PR1CS22}: $36\text{C} + 3\times30\text{C} = 132\text{C}$} \\
\multicolumn{4}{|l|}{\footnotesize \cite{FXFM23}: $5\times12\text{C} + 3\times4\text{C} + 4\text{C} = 76\text{C}$} \\
\multicolumn{4}{|l|}{\footnotesize Ours: $2\times12\text{C} + 3\times4\text{C} + 5\text{C} = 41\text{C}$\quad (excl.\ Fiat-Shamir overhead)} \\
\hline
\end{tabular}
\end{table}


\subsection{Operation Polynomials}\label{sec:op-polys}

We now state the verification equation for each step of Algorithm~\ref{alg:pairing} in a 3-pairing configuration (as used in Groth16 verification). All equations share the form $\text{LHS} = R(x) + Q(x) \cdot P_{12}(x)$; the quotient $Q$ is amortized across all steps via batching (Section~\ref{sec:iop:batching}).

\subsubsection{Zero-bit step}\label{sec:op:ml-0bit}

Lines~\ref{pr:ml:sq}--\ref{pr:ml:ldbl}: one squaring and three doubling lines.

\begin{equation}\label{eq:op:0b}
F^2(x) \cdot L_1(x) \cdot L_2(x) \cdot L_3(x) = R(x) + Q(x) \cdot P_{12}(x)
\end{equation}

$F, R \in \Fe{12}$ (12 coefficients each); $L_1, L_2, L_3 \in \text{Sparse}_{01379}$ (4 coefficients each). \textbf{Cost: 41C, 12 commitments.}

\subsubsection{Non-zero bit step}\label{sec:op:ml-nzbit}

Lines~\ref{pr:ml:rw1bit}--\ref{pr:ml:ladd-1bit}: squaring, residue witness, and both doubling and addition lines for each pair.

\begin{equation}\label{eq:op:n0b}
F^2(x) \cdot W^{\pm1}(x) \cdot L_{\text{d1}}(x) \cdot L_{\text{d2}}(x) \cdot L_{\text{d3}}(x) \cdot L_{\text{a1}}(x) \cdot L_{\text{a2}}(x) \cdot L_{\text{a3}}(x) = R(x) + Q(x) \cdot P_{12}(x)
\end{equation}

$W^{-1}$ is used for negative bits; the equation structure is identical. $W \in \Fe{12}$ encapsulates the equivalence class relationship from~\cite{OPP24}. \textbf{Cost: 69C, 12 commitments.}

\subsubsection{Frobenius correction step}\label{sec:op:ml-frob}

Lines~\ref{pr:ml:pi1}--\ref{pr:ml:pi2}: no squaring; two sparse correction lines per pair for the $\pi_p$ and $-\pi_p^2$ mappings.

\begin{equation}\label{eq:op:fm}
F(x) \cdot L_{1\pi}(x) \cdot L_{2\pi}(x) \cdot L_{3\pi}(x) \cdot L_{1\pi^2}(x) \cdot L_{2\pi^2}(x) \cdot L_{3\pi^2}(x) = R(x) + Q(x) \cdot P_{12}(x)
\end{equation}

\textbf{Cost: 55C, 12 commitments.}

\subsubsection{Final exponentiation step}\label{sec:op:finx}

Lines~\ref{pr:ml:27r}--\ref{pr:ml:fe}: cubic scaling and the Frobenius component ($q - q^2 + q^3$) of the final exponentiation, following~\cite{OPP24}. $R = 1$ for a valid pairing.

\begin{equation}\label{eq:op:fx}
F(x) \cdot c^{-q}(x) \cdot c^{q^2}(x) \cdot c^{-q^3}(x) \cdot W^w(x) = 1 + Q(x) \cdot P_{12}(x)
\end{equation}

$w \in \{0,1,2\}$ selects the cubic root of unity scaling; $W^w$ is sparse with 2 nonzero coefficients. The Frobenius images of $c$ are provided as prover inputs. \textbf{Cost: 56C, 12 commitments.}

\subsection{Security}\label{sec:op-sec}

Each verification equation takes the form $P = \prod_i A_i(x) - Q(x) \cdot P_{12}(x) - R(x)$ where $\deg(P) < 2^8$.\footnote{For $k$-pair pairings, $\deg(P) \approx 88 + 18(k-3)$, well below $2^8$ for all practical $k$. This bound can be extended to $2^{10}$ with negligible soundness impact.} By the Schwartz-Zippel lemma (Section~\ref{sec:bg:sz}):

\textbf{Soundness:} $\delta_s \leq d/|\Fq| < 2^8/2^{254} = 2^{-246}$ for BN254.

\textbf{Completeness:} An honest prover always produces $P(x) = 0$ identically, so the verifier always accepts. Euclidean division guarantees unique $Q, R$ with $\deg(R) < \deg(P_{12}) = 12$.
\section{Interactive Oracle Proof}

Now that we have covered the core construction behind our contribution, let's take a holistic view of verification of a pairing computation with an interactive proof system. While the construction above works beautifully on it's own, we can further optimize the protocol by batching all polynomial verifications into a two round interactive proof.

\subsection{Batching Polynomial Identity Testing}

The whole process of polynomial verification can be batched into a single polynomial verification. Here's how a two round interactive proof system would look like,

The prover can commit all $R$s for all $n$ equations involved in pairing computation in the first interaction and receive a random challenge $z$. The challenge $z$ can then be used to batch all the $Q$s into a single polynomial as a consecutive powers random linear combination like this,

$$Q_{acc} = \sum_{i=0}^{n - 1}z^i \cdot Q_i(z)$$

The verifier then prepares a new challenge $c$ and checks,

$$\sum_{i=0}^{n - 1}z^i \cdot A_i(c) \cdot B_i(c) \cdots X_i(c) = \sum_{i=0}^{n-1}{z_i \cdot R_i(c) + Q_{acc}(c) \cdot P_{12}(c)}$$

or,

\begin{equation}
    \label{eq:iop-batched}
    \sum_{i=0}^{n - 1}z^i \cdot [A_i(c) \cdot B_i(c) \cdots X_i(c) - R_i(c)] = Q_{acc}(c) \cdot P_{12}(c)
\end{equation}

This reduces $n$ polynomial evaluations of degree ranging from 38 to 76 (as described in Tables \ref{tbl:op:ml-0bit}, \ref{tbl:op:ml-nzbit}, \ref{tbl:op:ml-frob}, \ref{tbl:op:finx}) to a single polynomial of degree 76.

\subsection{Prover}

Prover receives $P$ and $Q$ point pairs, residue witness as computed in \cite{OPP24} and internal input $W$ for 27th root of unity. Prover then computes the $R_i$ and $Q_i$.

\begin{algorithm}
\caption{PolyMulDivQR: Polynomial multiplication and Euclidean division}
\label{alg:prover_polydiv}
\begin{algorithmic}[1]
\Require $P_1, P_2, \ldots, P_n \in \Fq[x]$
\Statex Internal: Divisor $P_{12} \in \Fq[x]$
\Ensure $Q, R \in \Fq[x]$ where $\deg(R) < \deg(P_{12})$ and $\prod_{i=1}^{n} P_i = Q \cdot P_{12} + R$
\State $Q \leftarrow 0$, $R \leftarrow P_1$ \gComm{Initialize quotient and remainder}
\For{$i = 2$ to $n$}
    \State $R \leftarrow R \cdot P_i$ \gComm{Multiply polynomials into R}
\EndFor
\State $d_P \leftarrow \deg(P_{12})$, $d_R \leftarrow \deg(R)$, $l_P \leftarrow \text{leading\_coeff}(P_{12})$
\While{$d_R \geq d_{P}$}
\State $t \leftarrow \text{leading\_coeff}(R)/l_P \cdot x^{d_R - d_{P}}$
\State $Q \leftarrow Q + t$, $R \leftarrow R - t \cdot P_{12}$, $d_R \leftarrow \deg(R)$
\EndWhile
\State \Return $(Q, R)$ \gComm{Quotient and Remainder with $\deg(R) < \deg(P_{12})$}
\end{algorithmic}
\end{algorithm}

\begin{algorithm}
\caption{Pairing Prover: Preparing polynomials}
\label{alg:ppp:2r:prover}
\begin{algorithmic}[1]
\Require $P \in \G_1, Q \in \G_2, c \in \Fe{12}, w \in \{0,1,2\}$
\Statex Internal: $W, W^2 \in \Fe{12}$ \Comment{$W$ is 27th root of unity}
\Ensure Arrays $Q_s, R_s$ of quotients and remainders
\color{gray}
\Statex \small{Every operation involving $T, P \text{ or } Q$ is repeated for $T_i, P_i \text{ and } Q_i$ for multi-pairing.}
\color{black}
\Statex Write $s = 6t + 2 = \sum_{i=0}^{L-1} s_i 2^i$ where $s_i \in \{-1, 0, 1\}$ and $s_L = 1$
\State $c^{-1} \leftarrow 1 / c$ \Comment{Inverse Residue witness}
\State $f \leftarrow c^{-1}$ \Comment{Initialize with inverse residue witness}
\State $T \leftarrow Q$ \gComm{$T_i$ for each pair for multi-pairing}
\Statex
\For{$i = L - 2$ to $0$}
    \State $L_\text{T,T} \leftarrow l_{T,T}(P)$ \gComm{Repeat for $T_i$}
    \State $T \leftarrow [2]T$ \gComm{Repeat for $T_i$}
    \If{$s_i = 0$}
			\State $Q, R \leftarrow {\textit{PolyMulDivQR}}(f, f, L_\text{T,T})$ \Comment{Algorithm \ref{alg:prover_polydiv}}
    \ElsIf{$s_i = 1$}
				\State $Q, R \leftarrow {\textit{PolyMulDivQR}}(f, f, c^{-1}, L_\text{T,T}, l_{T,Q}(P))$ \Comment{Algorithm \ref{alg:prover_polydiv}}
        \State $T \leftarrow T + Q$ \gComm{Repeat for $T_i, Q_i$}
    \ElsIf{$s_i = -1$}
				\State $Q, R \leftarrow {\textit{PolyMulDivQR}}(f, f, c, L_\text{T,T}, l_{T,-Q}(P))$ \Comment{Algorithm \ref{alg:prover_polydiv}}
        \State $T \leftarrow T - Q$ \gComm{Repeat for $T_i, -Q_i$}
    \EndIf
    \State $f \leftarrow R$, $Q_s$.\textit{append}($Q$), $R_s$.\textit{append}($R$)
\EndFor
\Statex
\State $Q_1 \leftarrow \pi_p(Q), Q_2 \leftarrow \pi_p^2(Q)$ \gComm{Repeat for $Q'_i, T_i$}
\State $L_{T,Q_1} \leftarrow l_{T,Q_1}(P)$
\State $T \leftarrow T + Q_1$  \gComm{Repeat for $Q_{1i}, T_i$}
\State $Q, R \leftarrow {\textit{PolyMulDivQR}}(f, L_{T,Q_1}, l_{T,-Q_2}(P))$ \Comment{Algorithm \ref{alg:prover_polydiv}}
\State $f \leftarrow R$, $Q_s$.\textit{append}($Q$), $R_s$.\textit{append}($R$)

\Statex

\If{$w = 1$}
    \State $Q, R \leftarrow {\textit{PolyMulDivQR}}(f, W, c^{-q}, c^{q^2}, c^{-q^3})$ \Comment{Algorithm \ref{alg:prover_polydiv}}
\ElsIf{$w = 2$}
    \State $Q, R \leftarrow {\textit{PolyMulDivQR}}(f, W^2, c^{-q}, c^{q^2}, c^{-q^3})$ \Comment{Algorithm \ref{alg:prover_polydiv}}
\EndIf
\State $Q_s$.\textit{append}($Q$), $R_s$.\textit{append}($R$) \gComm{Final $Q_s$ and $R_s$ accumulation, skip $f$}

\Statex

\State \Return $Q_s, R_s$
\end{algorithmic}
\end{algorithm}

\subsection{Verifier}

Verifier receives all the remainders, $R_i$, and RLC of quotients - $\sum_{i=1}^{n}c_i  \cdot Q_i$.
The verifier then prepares all $c_i$ solely from the remainder coefficients. For the evaluation point $x$, she includes all coefficients from the quotient RLC in the Fiat-Shamir hash as well. This guarantees soundness in Fiat Shamir. Then she uses generated $c_i$ for aggregating $\text{LHS}_i$ evaluated at $x$. In the end, she subtracts the evaluation of quotients RLC at $x$ times $P_{12}(x)$ from the aggregated LHS evaluation. If the result is zero, according to Schwartz--Zippel lemma, all the provided $R_i$ are correct with a very high\footnote{$Pr[P(r_0, r_1, r_2,..., r_d) = 0] \leq \frac{d}{|\F|}$} probability.

\subsection{Performance}

Batching operations in Miller loop steps yielded substantial performance improvements. Processing all operations within a Miller loop step as a single polynomial and performing Schwartz-Zippel verification of polynomial multiplications significantly reduced the overall computational requirements of the pairing verification process.
Combined with residue witness technique for final exponentiation elimination\cite{OPP24}, these optimizations enabled processing of ZK proofs on Starknet even before the modulo builtin\cite{sn_mod_builtin} became available, which later greatly improved finite field arithmetic operations. These techniques were subsequently integrated into Garaga, which remains the most efficient pairing implementation by operation count on Starknet to date.
% https://github.com/shramee/cairo_pairing/blob/main/packages/bn254_u256/src/pairing/schzip/steps.cairo

\subsection{[wip] Protocol Overview}

\subsection{[wip] Single-Round Protocol}

\subsubsection{[wip] Protocol description (Algorithm)}
\subsubsection{[wip] Complexity analysis}


\subsection{Two-Round Protocol with RLC Batching}

\subsubsection{[wip] Protocol description (Algorithm)}
\subsubsection{[wip] Complexity analysis}
\subsubsection{[wip] Comparison with single-round}


\subsection{Three-Round Protocol with Native Field Optimization}

\subsubsection{[wip] Protocol description (Algorithm)}
\subsubsection{[wip] Complexity analysis}
\subsubsection{[wip] Architecture-specific benefits}


\bibliographystyle{alpha}
\bibliography{references}

\end{document}